% See http://texblog.org/2013/11/11/latexs-alternative-letter-class-newlfm/
% and http://www.ctan.org/tex-archive/macros/latex/contrib/newlfm
% for more information.
%
\documentclass[11pt,stdletter,orderfromtodate,sigleft]{newlfm}
\usepackage{blindtext, xfrac, animate, hyperref, pxfonts}

  \setlength{\voffset}{-0.45in}

\newlfmP{dateskipbefore=0pt}
\newlfmP{sigsize=14pt}
\newlfmP{sigskipbefore=6pt}

\newlfmP{Headlinewd=0pt,Footlinewd=0pt}

\namefrom{\vspace{-0.3in}Jeremy R. Manning}
\addrfrom{
	Dartmouth College\\
    Department of Psychological \& Brain Sciences\\
    HB 6207 Moore Hall\\
	Hanover, NH  03755}

\addrto{}
\dateset{\today}

\greetto{To the editors of \textit{PLOS ONE}:}



\closeline{Sincerely,}

\begin{document}
\begin{newlfm}

I have enclosed our manuscript entitled \textit{A Stylometric Application of
Large Language Models} to be considered for publication as a \textit{Research
Article}. The manuscript shows that large language models can be used to
distinguish the writings of different authors: an individual GPT-2 model,
trained from scratch on the works of a single author, predicts that author's
held-out text more accurately than text written by others.

Our central claim is supported by a systematic and fully reproducible set of
experiments. We trained individual models on the works of eight classic authors
and used cross-entropy loss as a measure of stylistic similarity, an approach we
term predictive comparison. To confirm that our findings are robust rather than
artifacts of a single training run, we trained models across multiple random
seeds and several text representations (full text, content words, function
words, and part-of-speech sequences). These ablations show that both content
words and function words contribute to author-specific signatures, whereas
grammatical structure alone is less distinctive. We further apply the approach
to a real attribution problem, supporting R. P. Thompson's authorship of the
well-studied 15\textsuperscript{th} book of the \textit{Oz} series, originally
attributed to L. F. Baum.

We designed the study with reproducibility as a priority, in keeping with the
standards of \textit{PLOS ONE}: all code and data needed to reproduce every
result, figure, and statistical test are openly available at
\url{https://github.com/ContextLab/llm-stylometry}. We believe the work will
interest the journal's broad, interdisciplinary readership, including
researchers in computational linguistics, digital humanities, and machine
learning.

This manuscript reports original primary research that is not under
consideration for publication elsewhere. An earlier version has been posted as a
preprint on arXiv (arXiv:2510.21958), consistent with \textit{PLOS ONE}'s
preprint policy. Thank you for considering our manuscript; I hope you will find
it suitable for publication in \textit{PLOS ONE}.


\end{newlfm}
\end{document}
